Random projections reduce dimension without looking at the data, and Achlioptas's sparse matrices make them cheap. In a coursework comparison of a dense $\pm 1$ matrix (Eq.~3) and a sparse one (Eq.~4, with two thirds of the entries zero) on five small data sets from the University of California, Irvine (UCI) repository, Eq.~4 outperformed on three of five data sets from one matrix draw and one split, with the best kernel chosen on the test set. Over 100 draws of the same pipeline, that conclusion is supported about half of the time, but the mean difference does not hold it up: Eq.~4 minus Eq.~3 is \N{rep:iris:best_eq4_minus_best_eq3} points on Iris and \N{rep:seeds:best_eq4_minus_best_eq3} on Seeds (unadjusted intervals, one split), and in the tuned study, no comparison of Eq.~4 with Eq.~3 survives Holm's correction. The tuned study was performed with 100 paired draws per setting, five matrix families, target sizes 1 to 24, tuned kernels, and principal component analysis (PCA) and all-feature baselines. Its findings: Sparse matrices (Eq.~4 and, at these small input dimensions, the very sparse matrices of Li, Hastie, and Church) stay within about a point of a dense Gaussian matrix in paired differences averaged over target sizes and then equally over the data sets other than Iris on which each family was run (balanced accuracy for Glass), with both the pooling and the exclusion of Iris chosen after the per-cell results were seen; with Iris, Eq.~3 and the very sparse matrices still trail by under a point, but with $p < 0.05$, and Eq.~4 by \N{mag:pool:eq4:gauss}. Across all five data sets and the four non-Gaussian families, though, only \N{z:equiv_vs_gauss} of the \N{z:n_vs_gauss} primary comparisons with the Gaussian show equivalence within $\pm 2$ points, mostly at larger $k$. On four-feature Iris, where Eq.~4 often leaves a whole column of the matrix zero, it trails the Gaussian by \N{mag:pd:iris:3:eq4:gauss} to \N{mag:pd:iris:2:eq4:gauss} points at $k = 1$ to 3, and only $k = 2$ survives Holm's correction (paired $t$-tests at the 0.05 level; a Wilcoxon test gives Holm-adjusted $p = \N{wxh:iris:2:eq4:gauss}$). One nonzero per column loses at every target size, by up to \N{z:opc_holm_max} points. PCA beats the Gaussian projection by \N{pda:iris:2:gauss:pca} to \N{pda:wine:2:gauss:pca} points at $k = 2$. The variance of the squared-distance ratio, averaged over pairs, matches its exact formula to within about \N{vm:max_dev_pct}\%, but apart from one nonzero per column, which has both the highest variance and the lowest accuracy, its order does not explain the accuracy order; no target size here meets the Johnson--Lindenstrauss sufficient condition.
